\subsection{Definition of Riesz spaces}

Recall that, on a set E, a vector space structure over the field R and an order structure are said to be compatible if they satisfy the following two axioms:

\((\mathrm{OVS}_{\mathrm{I}})\) The relation \(x \leqslant y\) implies \(x + z \leqslant y + z\) for all \(z \in \mathbb{E}\) .

\((\mathrm{OVS}_{\mathrm{II}})\) The relation \(x \geqslant 0\) implies \(\lambda x \geqslant 0\) for every scalar \(\lambda > 0\) .

The space E, equipped with these two structures, is called an ordered vector space (TVS, II, §2, No.~5).

Axiom (OVS \(_{I}\) ) signifies that the order structure and the additive group structure on E are compatible, in other words that E, equipped with these two structures, is an ordered group (A, VI, §1, No.~1).

Axiom (OVS \(_{\text{I}}\) ) implies that the relations \(x \leqslant y\) and \(x + z \leqslant y + z\) are equivalent. Similarly, it follows from (OVS \(_{\text{II}}\) ) that, for every scalar \(\lambda > 0\) , the relations \(x \leqslant y\) and \(\lambda x \leqslant \lambda y\) are equivalent, because \(\lambda^{-1} > 0\) and so the relation \(\lambda x \leqslant \lambda y\) implies \(\lambda^{-1}(\lambda x) \leqslant \lambda^{-1}(\lambda y)\) . One can therefore say that, in an ordered vector space, the translations and the homotheties of ratio \(> 0\) are automorphisms of the order structure; this fact is also expressed by saying that the order is invariant under every translation and every homothety of ratio \(> 0\) . Moreover, the symmetry \(x \mapsto -x\) is an isomorphism of the order structure of E onto the opposite order structure.

DEFINITION 1. --- An ordered vector space is said to be a Riesz space (or lattice-ordered vector space) \(^{1}\) if its order structure is a lattice-ordering (that is, if every pair of elements x, y of E has a supremum \(\sup(x,y)\) and an infimum \(\inf(x,y)\) ). \(^{2}\)

Example. The space \(\mathbf{R}^{\mathrm{A}}\) of all real-valued functions defined on any set A is a Riesz space (for the order relation \(\ll x(t) \leqslant y(t)\) for all \(t \in \mathrm{A} \gg\) ); for, any two real-valued functions \(x, y\) defined on A have a supremum (resp. an infimum) equal to the mapping \(t \mapsto \sup(x(t), y(t))\) (resp. \(t \mapsto \inf(x(t), y(t))\) ).

One can also say that a Riesz space is a vector space E equipped with an order structure such that, on the one hand, this structure and the additive group structure of E define a lattice-ordered group structure on E (A, VI, §1, No.~9), and on the other hand that the axiom (OVS \(_{II}\) ) is satisfied.

Thus, all of the properties of lattice-ordered groups are applicable to Riesz spaces; we shall review here the principal ones (cf.~A, VI, §1, Nos. 9 to 12), as well as indicating the consequences that flow from the axiom (OVS \(_{II}\) ).

We recall first that one writes \(x^{+} = \sup(x,0)\) (the positive part of x), \(x^{-} = (-x)^{+} = \sup(-x,0)\) (the negative part of x), \(|x| = \sup(x,-x)\) (the absolute value of x); then \(x = x^{+} - x^{-}\) and \(|x| = x^{+} + x^{-}\) ; here, these two relations are equivalent to

\[
x ^ {+} = \frac {1}{2} (| x | + x), \quad x ^ {-} = \frac {1}{2} (| x | - x).
\]

The relation \(x \leqslant y\) is equivalent to \(\ll x^{+} \leqslant y^{+}\) and \(x^{-} \geqslant y^{-}\) . For any \(x\) and \(y\) , the triangle inequality holds:

\[
\vert x + y \vert \leqslant \vert x \vert + \vert y \vert .\tag{1}
\]

By the invariance of order under every homothety of ratio \textgreater{} 0,

\[
\sup (\lambda x, \lambda y) = \lambda \sup (x, y) \quad \text { for   all } \lambda \geqslant 0.\tag{2}
\]

In particular,

\[
(\lambda x) ^ {+} = \lambda x ^ {+}, (\lambda x) ^ {-} = \lambda x ^ {-} \quad \text { for   all } \lambda \geqslant 0.\tag{3}
\]

On the other hand, for \(\lambda < 0\) we have \((\lambda x)^{+} = (-\lambda x)^{-} = |\lambda |x^{-}\) and \((\lambda x)^{-} = (-\lambda x)^{+} = |\lambda |x^{+}\) ; it follows that, for every \(\lambda \in \mathbf{R}\) and every \(x\in \mathrm{E}\) ,

\[
\left| \lambda x \right| = \left| \lambda \right| \cdot \left| x \right|.\tag{4}
\]

The invariance of order under translation shows that for all \(z \in E\) ,

\[
\sup (x + z, y + z) = z + \sup (x, y),\tag{5}
\]

whence, in particular,

\[
\sup (x, y) = x + (y - x) ^ {+} = \frac {1}{2} (x + y + | x - y |).\tag{6}
\]

We have the relations

\[
\inf (x, y) = - \sup (- x, - y),\tag{7}
\]

\[
\sup (x, y) + \inf (x, y) = x + y.\tag{8}
\]

If \(x, y, z\) are \(\geqslant 0\) then (A, VI, §1, No.~12, Prop. 11)

\[
\inf (x + y, z) \leqslant \inf (x, z) + \inf (y, z).\tag{9}
\]

If A and B are two subsets of E each of which has a supremum, then \(A + B\) also has a supremum and

\[
\sup (A + B) = \sup A + \sup B.\tag{10}
\]

Two elements \(x, y\) of \(\mathrm{E}\) are said to be alien \(^3\) (to each other) if \(\inf(|x|, |y|) = 0\) ; by (8), this relation is equivalent to \(\sup(|x|, |y|) = |x| + |y|\) , and also, by (6), to \(||x| - |y|| = |x| + |y|\) ; 0 is the only element alien to itself; for every \(x \in \mathrm{E}\) , \(x^+\) and \(x^-\) are alien and may be characterized as the only alien elements \(y \geqslant 0\) , \(z \geqslant 0\) such that \(x = y - z\) . If \(y\) is alien to \(x\) , then every \(z \in \mathrm{E}\) such that \(|z| \leqslant |y|\) is also alien to \(x\) . If \(y\) and \(z\) are alien to \(x\) , then so is \(|y| + |z|\) , by the inequality (9); in particular, \(n|y|\) is alien to \(x\) for every integer \(n > 0\) , from which it follows that \(\lambda y\) is alien to \(x\) for every scalar \(\lambda\) , since there exists an integer \(n\) such that \(|\lambda| \leqslant n\) , whence \(|\lambda y| \leqslant n|y|\) . If a subset \(A\) of \(E\) consists of elements alien to \(x\) and if \(A\) has a supremum, then that supremum is also alien to \(x\) (A, VI, §1, No.~12, Cor. of Prop. 13).

Finally, we have the decomposition lemma (A, VI, §1, No.~10, Th. 1):

If \((x_{i})_{i\in \mathrm{I}}\), \((y_{j})_{j\in \mathrm{J}}\) are two finite sequences of elements \(\geqslant 0\) of E such that \(\sum_{i\in \mathrm{I}}x_i = \sum_{j\in \mathrm{J}}y_j\) , then there exists a finite sequence \((z_{ij})_{(i,j)\in \mathrm{I}\times \mathrm{J}}\) of elements \(\geqslant 0\) of E such that \(x_{i} = \sum_{j\in \mathrm{J}}z_{ij}\) for all \(i\in \mathrm{I}\) , and \(y_{j} = \sum_{i\in \mathrm{I}}z_{ij}\) for all \(j\in \mathrm{J}\) .

